%قبل از شروع مسئله ترجیحاً چیزی ننویسید. اما اگر خواستید بنویسید، در قالب لیستی (فقرات و شمارش) بنویسید.
%برای اضافه کردن نمره، بعدی از شروع‌مسئله، [x] را اضافه کنید و به جای x  نمره مورد نظر را قرار دهید. برای زیرمسئله‌ها هم به همین صورت است.
%\شروع‌مسئله[10]
\شروع‌مسئله

% صورت اصلی مسئله را در این‌جا بنویسید.
در کلاس درس، با الگوریتم
\lr{SELECT}
آشنا شدید که در مرتبه‌ی زمانی خطی، 
$k$ 
امین کوچک‌ترین عنصر آرایه‌ای نامرتب را به دست می‌داد. در عملکرد این الگوریتم مشاهده کردید که عناصر آرایه به دسته‌های پنج‌تایی تقسیم می‌شوند. حال اگر در الگوریتمی مشابه، عناصر را به دسته‌های هفت‌تایی تقسیم کنیم و سایر بخش‌های الگوریتم مانند قبل باشد:

\begin{enumerate}[ ]
    \item 
    حداقل فاصله‌ی میانه‌ی میانه‌ها از دو سر آرایه چقدر خواهد بود؟
 
    \item 
    فرض کنید برای مرتب‌سازی دسته‌ها از مرتب‌سازی حبابی
\lr{(Bubble sort)}
استفاده می‌شود. رابطه‌ی بازگشتی الگوریتم جدید را بنویسید و علت وجود هر جمله‌ی آن را به اختصار شرح دهید.

\item
نشان دهید الگوریتم همچنان عملکرد زمانی خطی دارد.

\item
اگر به جای دسته‌های پنج‌تایی یا هفت‌تایی از دسته‌های سه‌تایی استفاده شود چطور؟ مرتبه‌ی زمانی الگوریتم 
همچنان خطی خواهد بود؟
\end{enumerate}
 
\solution{ 
\begin{enumerate}[]
 \item 
    عناصر میانی نصف دسته‌ها از میانه‌ی میانه‌ها کوچک‌تر هستند. پس:
\begin{equation*}
\frac{1}{2} \times 4 \times \frac{n}{7} = \frac{2n}{7}
\end{equation*}

حداقل فاصله‌ی میانه‌ی میانه‌ها از دو سر آرایه
$\frac{2n}{7}$
می‌باشد.

\item 
 
مراحل الگوریتم و مرتبه‌ی زمانی مورد نیاز هر مرحله به قرار زیر است:
\begin{enumerate}
\item
دسته‌بندی آرایه‌ی اصلی به گروه‌های هفت‌تایی: O(n)
\item 
مرتب‌سازی هر دسته (در بدترین حالت نیاز به $7 \times 6$ مقایسه در هر دسته داریم) که نتیجه می‌دهد تعداد کل مقایسه‌ها:
\begin{equation*}
\frac{n}{7} \times 7 \times 6 \times C = 6Cn = O(n)
\end{equation*}\item
تشکیل آرایه‌ای از عناصر میانی دسته‌ها: $O(n)$
\item
 پیدا کردن بازگشتی میانه‌ی‌میانه‌ها: 
 $T(\frac{n}{7})$

\item
پارتیشن‌بندی آرایه‌ی اصلی با استفاده از میانه‌ی میانه‌ها به عنوان محور: $O(n)$
\item
پیداکردن k امین کوچک‌ترین عنصر به طور بازگشتی: 
$T(\frac{5n}{7})$

در نتیجه‌ رابطه‌ی بازگشتی به صورت زیر نوشته می‌شود:
\begin{equation*}
T(n) = T(\frac{n}{7}) + T(\frac{5n}{7}) + O(n)
\end{equation*}
\end{enumerate}

\item
 
حدس:
\begin{align*}
T(1)=c_1, O(n)=c_2n \hspace{10mm}
T(n) \leq 10(c_1 + c_2)n \\
\rightarrow T(n) \leq (\frac{6}{7})(10(c_1+c_2)) + c_2n = (\frac{60}{7})c_1n + (\frac{67}{7})c_2n \leq 10(c_1+c_2)n \\
\rightarrow T(n) = O(n)
\end{align*}



\item

در آن صورت:
\begin{equation*}
T(n)=T(\ceil{\frac{n}{3}})+T(\frac{4n}{6})+O(n) \geqslant T(\frac{n}{3})+T(\frac{2n}{3})+O(n)
\end{equation*}
حدس:
\begin{equation*}
T(m)=T(\frac{m}{3})+T(\frac{2m}{3})+O(m) \rightarrow 
T(m) \geqslant cm \log{m}  
\end{equation*}
\begin{align*}
\rightarrow T(n) \geqslant c(\frac{n}{3}) \log{(\frac{n}{3})}+c(\frac{2n}{3}) \log{(\frac{2n}{3})} + O(n) \\
\geqslant cn \log{n} + O(n)
\end{align*}پس در این حالت مرتبه‌ی زمانی خطی ندارد.

\end{enumerate}

}



%%%%%%%%%%%%%%%%%%%%%%%%%%%%%%%%%%%%%%%%%%%%%%%%%%%%%
% برخی دستورات مفید برای تمرین‌ها. %
%%%%%%%%%%%%%%%%%%%%%%%%%%%%%%%%%%%%%%%%%%%%%%%%%%%%%

% فرمول‌های نمونه:
%$f(x) = \Theta(x)$
%$T(x, y) \leq \log(x) + 2^y + \binom{2n}{n}$
% {\tt code\_function}


% لیست موارد بدون شماره‌گذاری:
%\شروع{فقرات}
%    \فقره مورد اول 
%        \شروع{فقرات}
%            \فقره مورد اول داخلی 
%                \شروع{فقرات}
%                    \فقره اولین مورد خیلی داخلی 
%                    \فقره دومین مورد خیلی داخلی 
%                \پایان{فقرات}
%            \فقره مورد دوم داخلی 
%        \پایان{فقرات}
%    \فقره مورد دوم 
%\پایان{فقرات}

% لیست موارد شماره‌گذاری شده:
%\شروع{شمارش}
%    \فقره یک مورد 
%    \فقره یک مورد دیگر 
%    \فقره باز هم یک مورد دیگر
%\پایان{شمارش}

% برای نوشتن یک فرمول چند خطی که یک شماره بخوره. برای حالت‌های دیگر به سایت زیر مراجعه کنید:
% https://www.overleaf.com/learn/latex/Aligning_equations_with_amsmath
%\begin{align} 
%    \begin{split}
%        (x+y)^3 &= (x+y)^2(x+y) \\
%                &= (x^2+2xy+y^2)(x+y) \\
%                &= (x^3+2x^2y+xy^2) + (x^2y+2xy^2+y^3) \\
%                &= x^3+3x^2y+3xy^2+y^3
%    \end{split}                                 
%\end{align}

% نوشتن ماتریس و آرایه دوبعدی
%\begin{align}
%    A = 
%    \begin{bmatrix}
%        A_{11} & A_{21} \\
%        A_{21} & A_{22}
%    \end{bmatrix}
%\end{align}

% اضافه کردن تصویر با شرح. برای اضافه کردن تصویر بدون شرح، فقط دستور درج‌تصویر را استفاده کنید.
%\شروع{شکل}[h!]
%\تنظیم‌ازوسط
%  \درج‌تصویر[پهنا=0.5\پهنای‌متن]{img.jpg}
%\شرح{شرح تصویر}
%\برچسب{تصویر:۱}
%\پایان{شکل}
%\پایان{زیرمسائل}

%کد زیر را ببینید.
%%\lstinputlisting{source_filename.py}
%\begin{latin}
%\begin{lstlisting}
%# Import the modules
%import sys
%import random
%
%ans = True
%
%while ans:
%question = raw_input("Ask the magic 8 ball a question: (press enter to quit) ")
%
%answers = random.randint(1,8)
%
%if question == "":
%sys.exit()
%
%elif answers == 1:
%print "It is certain"
%
%elif answers == 2:
%print "Outlook good"
%
%elif answers == 3:
%print "You may rely on it"
%
%elif answers == 4:
%print "Ask again later"
%
%elif answers == 5:
%print "Concentrate and ask again"
%
%elif answers == 6:
%print "Reply hazy, try again"
%
%elif answers == 7:
%print "My reply is no"
%
%elif answers == 8:
%print "My sources say no"
%\end{lstlisting}
%\end{latin}
